\textbf{Phần 1. Trắc nghiệm lựa chọn}
\setcounterpageref{ex}{0}

\begin{ex} %[2H6V30]
Một năm có bao nhiêu tháng?
\choice
{$10$ tháng}
{\True $12$ tháng}
{$11$ tháng}
{$14$ tháng}
\loigiai{
Một năm luôn có đúng 12 tháng.
\textbf{Đáp án của câu hỏi là B}
}
\end{ex}

\begin{ex} %[2H6V30]
Tháng nào dưới đây có $31$ ngày?
\choice
{\True Tháng 1}
{Tháng 4}
{Tháng 6}
{Tháng 9}
\loigiai{
Các tháng có 31 ngày là: Tháng 1, 3, 5, 7, 8, 10, 12.
\textbf{Đáp án của câu hỏi là A}
}
\end{ex}

\begin{ex} %[2H6V30]
Một tuần lễ có bao nhiêu ngày?
\choice
{$5$ ngày}
{$6$ ngày}
{\True $7$ ngày}
{$8$ ngày}
\loigiai{
Một tuần lễ có 7 ngày.
\textbf{Đáp án của câu hỏi là C}
}
\end{ex}

\begin{ex} %[2H6V30]
Nếu hôm nay là thứ Ba ngày $5$ tháng $10$. Thứ Ba tuần sau là ngày mấy tháng $10$?
\choice
{Ngày $10$}
{Ngày $11$}
{Ngày $13$}
{\True Ngày $12$}
\loigiai{
Thứ Ba tuần sau là: $5 + 7 = 12$.
\textbf{Đáp án của câu hỏi là D}
}
\end{ex}

\begin{ex} %[2H6V30]
Ngày Quốc tế Thiếu nhi $1$ tháng $6$ năm nay rơi vào thứ Sáu. Hỏi ngày $8$ tháng $6$ năm nay là thứ mấy?
\choice
{Thứ Bảy}
{\True Thứ Sáu}
{Thứ Năm}
{Chủ nhật}
\loigiai{
Từ ngày 1 đến ngày 8 là $8 - 1 = 7$ ngày (đúng 1 tuần). Do đó vẫn là thứ Sáu.
\textbf{Đáp án của câu hỏi là B}
}
\end{ex}

\begin{ex} %[2H6V30]
Tháng nào trong năm có ít ngày nhất?
\choice
{\True Tháng 2}
{Tháng 4}
{Tháng 6}
{Tháng 11}
\loigiai{
Tháng 2 có 28 hoặc 29 ngày, ít ngày nhất trong năm.
\textbf{Đáp án của câu hỏi là A}
}
\end{ex}

\begin{ex} %[2H6V30]
Nếu thứ Bảy tuần này là ngày $18$ thì thứ Bảy tuần trước là ngày mấy?
\choice
{Ngày $10$}
{Ngày $12$}
{\True Ngày $11$}
{Ngày $9$}
\loigiai{
Thứ Bảy tuần trước là: $18 - 7 = 11$.
\textbf{Đáp án của câu hỏi là C}
}
\end{ex}

\begin{ex} %[2H6V30]
Các tháng có $30$ ngày trong năm là
\choice
{Tháng 1, 3, 5, 7}
{Tháng 2, 4, 6, 8}
{Tháng 7, 8, 9, 10}
{\True Tháng 4, 6, 9, 11}
\loigiai{
Các tháng có 30 ngày gồm: Tháng 4, tháng 6, tháng 9 và tháng 11.
\textbf{Đáp án của câu hỏi là D}
}
\end{ex}

\begin{ex} %[2H6V30]
Nếu ngày $20$ của một tháng là Chủ nhật thì ngày Chủ nhật tiếp theo của tháng đó là ngày mấy?
\choice
{Ngày $26$}
{\True Ngày $27$}
{Ngày $28$}
{Ngày $25$}
\loigiai{
Chủ nhật tiếp theo là: $20 + 7 = 27$.
\textbf{Đáp án của câu hỏi là B}
}
\end{ex}

\begin{ex} %[2H6V30]
Một tháng có thể có nhiều nhất bao nhiêu ngày Chủ nhật?
\choice
{\True $5$ ngày Chủ nhật}
{$4$ ngày Chủ nhật}
{$6$ ngày Chủ nhật}
{$3$ ngày Chủ nhật}
\loigiai{
Một tháng có tối đa 31 ngày, $31 : 7 = 4$ dư 3 ngày. Do đó có tối đa 5 ngày Chủ nhật.
\textbf{Đáp án của câu hỏi là A}
}
\end{ex}

\begin{ex} %[2H6V30]
Bố đi công tác $2$ tuần lễ. Hỏi bố đi công tác bao nhiêu ngày?
\choice
{$10$ ngày}
{$12$ ngày}
{\True $14$ ngày}
{$16$ ngày}
\loigiai{
$2 \times 7 = 14$ (ngày).
\textbf{Đáp án của câu hỏi là C}
}
\end{ex}

\begin{ex} %[2H6V30]
Nếu ngày cuối cùng của tháng là ngày $31$ thì tháng đó KHÔNG THỂ là
\choice
{Tháng 7}
{Tháng 8}
{Tháng 12}
{\True Tháng 9}
\loigiai{
Tháng 9 chỉ có 30 ngày nên không thể kết thúc vào ngày 31.
\textbf{Đáp án của câu hỏi là D}
}
\end{ex}

\textbf{Phần 2. Trắc nghiệm đúng sai}
\setcounterpageref{ex}{0}

\begin{ex} %[2H6V30]
Xét tính đúng sai của các nhận định về ngày tháng trong năm:
\choiceTF
{\True Tháng 7 và Tháng 8 là hai tháng liên tiếp đều có 31 ngày.}
{Tháng 2 luôn luôn chỉ có 28 ngày.}
{\True Nếu ngày 2 là thứ Hai thì ngày 9 cùng tháng cũng là thứ Hai.}
{Tất cả các tháng chẵn trong năm đều có 30 ngày.}
\loigiai{
\begin{itemize}
    \item Ý a) Đúng: Đúng, cả tháng 7 và tháng 8 đều có 31 ngày.
    \item Ý b) Sai: Năm nhuận tháng 2 có 29 ngày.
    \item Ý c) Đúng: $2 + 7 = 9$, đúng 1 tuần sau.
    \item Ý d) Sai: Tháng 8, 10, 12 là tháng chẵn nhưng có 31 ngày; tháng 2 có 28/29 ngày.
\end{itemize}
\begin{itemize}
    \item a) \textbf{Đúng}.
    \item b) \textbf{Sai}.
    \item c) \textbf{Đúng}.
    \item d) \textbf{Sai}.
\end{itemize}
}
\end{ex}

\begin{ex} %[2H6V30]
Trong một tháng có ngày thứ Hai rơi vào các ngày: ngày 3, ngày 10, ngày 17, ngày 24, ngày 31:
\choiceTF
{\True Tháng đó có đúng 5 ngày thứ Hai.}
{\True Tháng đó chắc chắn phải có 31 ngày.}
{Ngày 1 của tháng đó là ngày Chủ nhật.}
{\True Thứ Ba đầu tiên của tháng đó là ngày 4.}
\loigiai{
\begin{itemize}
    \item Ý a) Đúng: Các ngày: 3, 10, 17, 24, 31 (5 ngày).
    \item Ý b) Đúng: Vì có ngày 31 là thứ Hai nên tháng đó có 31 ngày.
    \item Ý c) Sai: Ngày 3 là thứ Hai thì ngày 2 là Chủ nhật, ngày 1 là thứ Bảy.
    \item Ý d) Đúng: Sau ngày 3 là ngày 4 (thứ Ba). Đúng.
\end{itemize}
\begin{itemize}
    \item a) \textbf{Đúng}.
    \item b) \textbf{Đúng}.
    \item c) \textbf{Sai}.
    \item d) \textbf{Đúng}.
\end{itemize}
}
\end{ex}

\textbf{Phần 3. Trả lời ngắn}
\setcounterpageref{ex}{0}

\begin{ex} %[2H6V30]
Một tuần có 7 ngày. Hỏi 3 tuần lễ có tất cả bao nhiêu ngày?
\shortans{21}
\loigiai{
\textbf{Đáp số:} $21$\\
$7 \times 3 = 21$ ngày.
}
\end{ex}

\begin{ex} %[2H6V30]
Nếu thứ Năm tuần này là ngày 12 thì thứ Năm tuần sau là ngày mấy?
\shortans{19}
\loigiai{
\textbf{Đáp số:} $19$\\
$12 + 7 = 19$.
}
\end{ex}

\begin{ex} %[2H6V30]
Tháng 12 có bao nhiêu ngày?
\shortans{31}
\loigiai{
\textbf{Đáp số:} $31$\\
Tháng 12 có 31 ngày.
}
\end{ex}

\begin{ex} %[2H6V30]
Năm nay ngày 20 tháng 11 là ngày thứ Tư. Hỏi ngày 27 tháng 11 năm nay là thứ mấy? (Quy ước: Thứ Hai = 2, ..., Thứ Tư = 4, Chủ nhật = 8)
\shortans{4}
\loigiai{
\textbf{Đáp số:} $4$\\
Cách nhau 7 ngày nên vẫn là thứ Tư (số 4).
}
\end{ex}

\begin{ex} %[2H6V30]
Trong một năm có bao nhiêu tháng có 30 ngày?
\shortans{4}
\loigiai{
\textbf{Đáp số:} $4$\\
Có 4 tháng: tháng 4, 6, 9, 11.
}
\end{ex}

\begin{ex} %[2H6V30]
Kỳ nghỉ hè của học sinh kéo dài trong 8 tuần lễ. Hỏi kỳ nghỉ hè kéo dài bao nhiêu ngày?
\shortans{56}
\loigiai{
\textbf{Đáp số:} $56$\\
$8 \times 7 = 56$ ngày.
}
\end{ex}
